Feature attributions such as SHAP and LIME are routinely read as statements about which variables drive an outcome, although they describe a fitted model rather than the data-generating process, and in practice the causal truth needed to notice the difference is unavailable. We ask whether the error of such an explanation---its distance from the true total interventional effect---can be predicted from quantities that are observable without causal ground truth. We build a benchmark of 300 structural causal models from five graph families with exact causal effects and exact Shapley values of the true function, train XGBoost and MLP models with three seeds, run five explainers, and decompose each explanation's error into an estimand gap, model misfit and explainer approximation error, measured on importance shares and in rank space. The estimand gap dominates: with a perfect model and exact SHAP the importance shares are already $0.31$ ($L_2$) from the causal shares, model misfit adds $0.09$--$0.14$, and KernelSHAP's own error is below $0.01$; the gap depends on the causal structure and is largest for chain-mediation worlds, where 38\% of the causal effect sits on features that exact SHAP credits with nothing. From 39 ground-truth-free observables---dependence, model fit, attribution concentration, stability and disagreement---a gradient-boosted meta-model predicts the share error with Spearman $\rho = 0.65$ on unseen SCMs and $0.52$ (AUROC $0.76$) under leave-one-graph-family-out, beating every single-signal baseline ($p < 0.001$, paired bootstrap over SCMs). Predicting the components separately shows where this works: explainer error ($\rho = 0.91$) and model misfit ($0.70$) are visible in the observables, the structural estimand gap ($0.34$) is largely not, and the feature ranking error ($0.22$) is hardly predictable at all. A diagnostic can therefore warn when an explanation is unreliable because of the model or the explainer, but not when SHAP is answering a different question from the causal one---which is the dominant failure.